\documentclass[11pt]{article}

\usepackage[margin=1in]{geometry}
\usepackage[T1]{fontenc}
\usepackage{lmodern}
\usepackage{amsmath,amssymb}
\usepackage{enumitem}
\usepackage{hyperref}
\usepackage{xcolor}
\usepackage{parskip}

\hypersetup{
  colorlinks=true,
  linkcolor=blue!50!black,
  citecolor=blue!50!black,
  urlcolor=blue!50!black,
  pdftitle={Time Decay in Fellegi--Sunter Record Linkage: A Literature Survey},
  pdfauthor={Denis Robert},
}

\title{Time Decay in Fellegi--Sunter Record Linkage:\\A Literature Survey}
\author{Denis Robert}
\date{August 2026}

\emergencystretch=2em
\newcommand{\dt}{\Delta t}

\begin{document}
\maketitle

\begin{abstract}
This survey reviews the documented line of work that treats comparison evidence in
probabilistic (Fellegi--Sunter) record linkage as \emph{decaying} with the gap between two
records. The static $m/u$ parameters of the classical model have, since 2011, been extended
with temporal variants in which agreement probabilities, comparison weights, or the final
similarity score become functions of the elapsed time between record timestamps. We
summarize the primary points of the founding temporal-linkage papers and of the works closest
to the present project: Li et al.\ (2011), Christen and Gayler (2013), Chiang et al.\ (2014),
Li et al.\ (2015), Hu et al.\ (2017), Ranbaduge and Christen (2018/2020), Litoux and Ray
(2026), and Shim's \emph{TimeLink} (2026). We then present an information-theoretic reading:
a Fellegi--Sunter weight is a log-likelihood ratio measured in bits, its match-conditional
expectation is a Kullback--Leibler divergence, and $u$ is tied, in the frequency model, to a
R\'enyi collision entropy; ``decay of evidence with age'' is therefore the statement that the
mutual information between a comparison outcome and match status falls to zero over time. The
survey closes with a gap analysis of what, to our knowledge, remains unpublished.
\end{abstract}

\section{Where decay can enter Fellegi--Sunter}

A Fellegi--Sunter (FS) linkage scores a candidate pair of records $(a,b)$ by summing, over
comparison fields $c$ and comparison levels $\ell_c$, the log-likelihood ratio
\begin{equation}
  \label{eq:fs}
  w_c(\ell_c)\;=\;\log_2\frac{P(\ell_c \mid M)}{P(\ell_c \mid U)}
                \;=\;\log_2\frac{m_c(\ell_c)}{u_c(\ell_c)},
\end{equation}
where $M$ is ``same entity,'' $U$ is ``different entity,'' $m_c$ is the probability of
agreement given a true match, and $u_c$ is the probability of agreement given a non-match
(the coincidence rate). Posterior mass is monotone in the summed weights~\cite{fellegi1969,winkler1993}. FS assumes
conditional independence across fields and, crucially for this survey, that $m_c$ and $u_c$
are \emph{constants}, independent of any covariate of the pair.

The temporal literature inserts ``decay'' at four distinct points:
\begin{enumerate}[leftmargin=*]
  \item \textbf{Blocking and similarity.} The retrieval score is attenuated by the
        timestamp gap (Li et al.\ 2011; Litoux and Ray 2026; the present project's
        multi-view blocker).
  \item \textbf{Weight exponent.} The per-field contribution $\log(m_c/u_c)$ is multiplied
        by a decay factor of the gap.
  \item \textbf{Field probability $m_c$.} The agreement probability itself becomes a
        function $m_c(\dt)$, converging towards the coincidence baseline as the gap grows;
        equivalently $P(\gamma_c\mid M,\dt)\to P(\gamma_c\mid U)$. This is where
        ``decay $m$, keep $u$ fixed'' sits (TimeLink; the present project).
  \item \textbf{Value transition models.} $m$ and $u$ are replaced by a full distribution
        over what a value becomes after $\dt$ (Chiang et al.\ 2014; Li et al.\ 2015;
        TimeLink's semi-Markov models).
\end{enumerate}

\section{Primary sources, chronologically}

\subsection{Li, Dong, Maurino, Srivastava (2011): first explicit time decay}

Pei Li, Xin Luna Dong, Andrea Maurino, Divesh Srivastava.
\emph{Linking temporal records}. PVLDB 4(11):956--967, 2011.
DOI: \url{https://doi.org/10.14778/3402707.3402733}~\cite{li2011}.

This is the paper that defined temporal record linkage: records carry a time-stamp and the
task is to identify records that describe the same entity over time. Two mechanisms are
proposed:
\begin{itemize}[leftmargin=*]
  \item \emph{Time decay}: ``We apply time decay to capture the effect of elapsed time on
        value evolution'' --- the first explicit use of the term in the linkage
        literature.
  \item \emph{Clustering rather than pairwise decisions}: ``Instead of comparing each pair
        locally, we propose clustering methods that consider the order of records to make
        global decisions.''
\end{itemize}
Experiments show the algorithms ``significantly outperform traditional record linkage
techniques on various temporal data sets.'' This paper establishes the vocabulary
(``time decay'') and the two axes --- per-pair weight decay and global/orderings-aware
decisions.

\subsection{Christen and Gayler (2013): adaptive temporal entity resolution}

Peter Christen, Ross W.\ Gayler.
\emph{Adaptive temporal entity resolution on dynamic databases}.
PAKDD 2013 (LNAI 7819), pages 558--569.
DOI: \url{https://doi.org/10.1007/978-3-642-37456-2\_47}~\cite{cg2013}.

The first work, to our knowledge, to make the effective $m$-probabilities explicitly
gap-dependent and updated from data.
\begin{itemize}[leftmargin=*]
  \item Streaming setup: query records are matched against a growing database; a matched
        query is inserted back into the database (online, adaptive).
  \item For each attribute and time-lag bin the system maintains agree/disagree counts from
        accepted matches, yielding a data-driven agreement probability
        $P(A_{\dt}\mid S)$; the disagreement probability for non-matches is
        frequency-based ($1-P_j(v)$, with $P_j$ the value frequency).
  \item Each attribute similarity $s_j$ is rescaled by a time-dependent weight: an
        agreement is down-weighted to the extent that different entities plausibly share
        the value; a disagreement is down-weighted to the extent that same entities
        plausibly changed it.
  \item Models fitted ``adaptively'' (update per insertion) and ``non-adaptively''
        (frozen training-time counters). On synthetic data with noise the adaptive model
        wins; on the NC Voter database (about 2.4M records) the \emph{adaptive} variant
        underperformed its frozen baseline, because with few true matches per query the
        counters are noisy --- a cautionary result directly relevant to the calibration
        paradox seen in the white paper.
\end{itemize}

\subsection{Chiang, Doan, Naughton (2014): entity evolution}

Yue-Hsuan Chiang, AnHai Doan, Jeffrey F.\ Naughton.
\emph{Modeling entity evolution for temporal record matching}.
SIGMOD 2014, pages 1175--1186.
DOI: \url{https://doi.org/10.1145/2588555.2588560}~\cite{cd14}.

The paper recognizes that ``any such method relies at its heart on a model that captures
information about how the entities evolve''. It proposes:
\begin{itemize}[leftmargin=*]
  \item a more detailed model of the \emph{reappearance} probability of an attribute value
        over time, conditionally dependent on past values, rather than a simple
        change-or-not decay;
  \item matching over \emph{sets} of records, not pairs, for robustness;
  \item experimental evidence of improved accuracy and noise-resistance at minimal
        overhead.
\end{itemize}
This is the paper that turns ``decay'' into ``stochastic value evolution,'' a full
transition model; TimeLink's semi-Markov formulation is in this lineage.

\subsection{Li, Lee, Hsu, Tan (2015)}

Furong Li, Mong Li Lee, Wynne Hsu, Wang-Chiew Tan.
\emph{Linking temporal records for profiling entities}.
SIGMOD 2015, pages 593--605.
DOI: \url{https://doi.org/10.1145/2723372.2737789}~\cite{ll15}.

The abstract introduces a transition model that ``captures the probability that an entity may
change to a particular attribute value after some time period'' and combines it with
\emph{source quality/freshness} metrics: timestamps, the value-transition model, and source
timeliness are jointly exploited to link temporal records into the correct time period of
the entity profile. Each source's freshness acts as an additional, source-level
``$u$ moving with era,'' which connects to the open question in this survey
(Section~\ref{sec:gap}, item~3).

\subsection{Hu, Wang, Vatsalan, Christen (2017): regression classification}

Yichen Hu, Qing Wang, Dinusha Vatsalan, Peter Christen.
\emph{Improving temporal record linkage using regression classification}.
PAKDD 2017 (LNCS 10235), pages 561--573.
DOI: \url{https://doi.org/10.1007/978-3-319-57454-7\_44}~\cite{hu2017}.

Verified against the accepted text. Key points:
\begin{itemize}[leftmargin=*]
  \item The decay probabilities of Li et al.\ are replaced by a \emph{logistic
        regression} of the (time-gap) disagreement probability $\hat d_{\neq}(A,\dt)$,
        with support attributes (gender, age) and the gap as features; the regression
        replaces the (estimated) decay, while agreement decay $d_=$ is retained.
  \item The weight on attribute $A$ becomes
        \begin{equation}
          w_A(s_A,\dt) = 1 - s_A\,d_=(A,\dt) - (1-s_A)\, d_{\neq}(A,\dt),
          \label{eq:hu}
        \end{equation}
        with an \emph{impact-reduction} coefficient $\alpha$ that damps the temporal
        blow (as $\alpha\to\infty$ the score re-approaches a non-temporal scorer).
  \item Evaluation on NC Voter subsets (5k--100k): at full strength the pure decayed
        model is poor; a moderate $\alpha$ and a larger training pool give the temporal
        signal a small win (``Mixed''). The paper explicitly reports that without the
        damping the decay models ``degraded'' the linkage quality.
\end{itemize}
Relevance to the project: this is the strongest published statement that a decayed $m$ is not
free; without value-frequency-aware $u$ and a regularized (hard) gap effect it can pass
below the untrained static baseline on real data --- a result that the project's
calibration experiments reproduce.

\subsection{Ranbaduge and Christen (2018, 2020): temporal and privacy}

Thilina Ranbaduge and Peter Christen.
2018: \emph{Privacy-preserving temporal record linkage}. ICDM 2018, pages 377--386.
DOI: \url{https://doi.org/10.1109/ICDM.2018.00053}~\cite{rc2018}.
2020: \emph{A scalable privacy-preserving framework for temporal record linkage}.
Knowledge and Information Systems 62(1):45--78.
DOI: \url{https://doi.org/10.1007/s10115-019-01370-1}~\cite{rc2020}.

These works operate at the intersection of temporal linkage and privacy-preserving
record linkage (Bloom filters and homomorphic encryption), and use the same
phenomenon --- that entity attributes evolve over time --- to adjust an encoded
similarity, obtaining ``better linkage quality compared to non-temporal PPRL'' on data
sets of several million records. Their contribution is a protocol that safely uses the
value-change probability, i.e.\ they make use of a time-conditioned $m/u$ without
revealing the raw values.

\subsection{Litoux and Ray (2026): two kinds of decay on vessel data}

Victor Litoux, Cyril Ray.
\emph{Temporal record linkage using time decay models applied to vessel data}.
IEEE MDM 2026, pages 313--318.
DOI: \url{https://doi.org/10.1109/MDM71479.2026.00044}~\cite{lr2026}.

An applied paper that explicitly separates \emph{disagreement decay} (the probability that
an entity's value changes across the gap) from \emph{agreement decay} (the weight given to
the shared attributes), applied to Automatic Identification System (AIS) vessel
trajectories: common attributes anchor the linkage and the score is the decayed
combination of these two effects. This is the closest item to the present project's
pair-conditional decay framing, though on vessel data rather than person records.

\subsection{Shim (2026): TimeLink}

Kwan Bo Shim.
\emph{TimeLink: time-aware record linkage with semi-Markov dynamics}.
SSRN preprint 6468518, 2026.
DOI: \url{https://doi.org/10.2139/ssrn.6468518}~\cite{shim2026}.

The closest published match of the project's \emph{model}: make $m$ time-dependent, keep
$u$ static. Details, from the preprint's own text:
\begin{itemize}[leftmargin=*]
  \item Field taxonomy: \emph{stable} (DOB, sex; fixed $m$, $u$), \emph{dynamic}
        (low-cardinality: marital status, ZIP; semi-Markov/CTMC transition matrix), and
        \emph{semi-stable} (high-cardinality: last name, address; a Weibull survival
        $S(\dt)$ on the current value).
  \item The match probability for an attribute is (eq. 8 of the paper)
        \begin{equation}
          m^{*}(\dt) = S(\dt)\cdot m_{\mathrm{same}} + [1-S(\dt)]\cdot
          m_{\mathrm{changed}},\qquad m_{\mathrm{changed}} \approx u_c,
        \end{equation}
        which, for $S(\dt)=e^{-\dt/\bar\tau}$, is precisely
        $m(\dt)=u+(m-u)\,e^{-\dt/\bar\tau}$ --- the form already used in the project's
        design note.
  \item The weights are the FS form with the time-dependent $m$:
        $w_{\mathrm{agree}}=\ln(m(\dt)/u)$ and
        $w_{\mathrm{disagree}}=\ln\bigl((1-m(\dt))/(1-u)\bigr)$.
  \item Experiments: NC Voter (10-year gap): TimeLink F1 $0.877$ vs.\ Splink $0.751$;
        IPUMS historical census (20-year gap): $0.668$ vs.\ $0.146$; a synthetic 10--30
        year study shows classical FS collapsing to $0.055$ while TimeLink stays above
        $0.85$.
\end{itemize}
Positioning relative to TimeLink: both make $m$ decayable and $u$ static; TimeLink derives
$m(\dt)$ from generative sojourn-time models and precomputed lookup tables, while the
present project proposes an additive pair-conditional term on the logged $\log_2(m_c/u_c)$
stack that composes with any $m/u$ calibration, joined with measured per-field rates on a
FAISS/Splink pipeline --- a measurement, not just a model.

\section{The entropy reading of decay}

The literature as summarized converges on ``transient fields decay, stable fields do not''
but leaves the form and the benefit-conditional evidence open. The present section records an
information-theoretic view that unifies these observations.

\begin{itemize}[leftmargin=*]
  \item \textbf{FS weights are information in bits.} Each contribution in
        \eqref{eq:fs} is a log-likelihood-ratio: its match-conditional expectation is the
        Kullback--Leibler divergence
        \begin{equation}
          \mathbb{E}\!\left[\log_2 \frac{P(\gamma\mid M)}{P(\gamma\mid U)} \right]
          \;=\; \sum_\gamma P(\gamma\mid M)\log_2
            \frac{P(\gamma\mid M)}{P(\gamma\mid U)}
          \;=\; D_{\mathrm{KL}}(M_c \| U_c)
          \;\ge\;0,
        \end{equation}
        the expected ``surprise'' of seeing the pair (agree vs.\ disagree) in the match vs.\
        chance case.'' A decayed weight of 0 carries no information: the outcome is
        equally expected under $M$ and $U$.
  \item \textbf{$u_c$ is a collision (R\'enyi) entropy.} In a frequency model,
        $u_c(v) = f_c(v)$, the value frequency. The probability of a random coincidental
        agreement is the quadratic collision
        $\sum_v f(v)^2 = \exp(-H_2)$ with $H_2$ the R\'enyi-2 entropy of the value
        distribution. So ``common value vs.\ rare value'' is exactly a difference of
        $H_2$: a rare agreement carries more surprise --- $-\log_2 u$ worth --- the
        entropy content of the field.
  \item \textbf{Decay erases mutual information.} Under the time-dependent model
        $P(\cdot\mid M,\dt)\to P(\cdot\mid U)$, the KL diverges to zero and the pointwise
        weights go to zero; the mutual information
        $I(\gamma; M \mid \dt)$ between the comparison outcome and the match
        indicator falls to zero as the gap ages. For the exponential
        $m(\dt)=u+(m-u)e^{-\dt/T}$, the pointwise weight halves after about $T\ln 2$.
  \item \textbf{Noise entropy, not loss of the data.} The decay raises the
        match conditional entropy $H(P(\cdot \mid M,\dt))$ towards
        $H(P(\cdot \mid U))$: the observed contact fields of a true pair become as
        uninformative as the fields of an arbitrary pair. Decay therefore is a statement
        about the \emph{expected entropy of the evidentiary outcome}, not about the
        information being lost from the data set --- the surname is still itself; it is
        just no longer evidence.
\end{itemize}

\noindent Crucially, entropy is concerned with \emph{how much} information, not its
\emph{direction}; the caveats recorded for the temporal literature above are not
alternatives to this reading. The entropy framing explains the sign of a known
phenomenon --- why pairs at long gaps should trust identity anchors and discount transient
contact fields --- and provides per-field ``half-life'' numbers that are directly
measurable from a labelled corpus, which is what the measured experiments in the present
project attempt.

\section{Against the ``ad hoc decay function'' objection}

TimeLink's central methodological complaint about earlier temporal work is that it
``relies on ad hoc decay functions rather than deriving weights from a coherent
generative model.'' The entropy reading suggests that the decay is not a free knob to
be pulled but the \emph{only} information-theoretically justified completion of the
static FS model once a field's lifetime enters. This section makes that precise.

\paragraph{Step 1: the weight keeps its likelihood-ratio form.}
A Fellegi--Sunter weight is a Bayes factor, and Bayes factors are not chosen: they are
determined by the model for the comparison outcome. Adding the capture-date gap $\dt$
as a covariate does not change the decision rule --- the optimal evidence for deciding
$M$ versus $U$ given $\gamma$ and $\dt$ is always the ratio
\begin{equation}
  w_c(\ell;\dt) \;=\; \log_2\frac{P(\ell \mid M,\dt)}{P(\ell \mid U,\dt)}.
  \label{eq:llr-time}
\end{equation}
There is no alternative weight family: any other scalar evidence leaves
discrimination bits on the table (Neyman--Pearson optimality). So ``decay'' cannot be
introduced as a new modeling act; it is forced once the model conditions on $\dt$.

\paragraph{Step 2: decompose into a baseline and a max-entropy survival part.}
Write the match-conditional agreement probability as
\begin{equation}
  m_c(\dt) \;=\; u_c + \bigl(m_c(0)-u_c\bigr)\, S_c(\dt),
  \label{eq:surv}
\end{equation}
where $S_c(\dt)\in[0,1]$, $S_c(0)=1$, is the survival of the field's current value, and
$m_c(0)$ is its modern agreement probability. The static model is the special case
$S_c\equiv1$. Which functional form for $S_c$? The information-theoretic default is
\emph{maximal ignorance}: given only the mean residency $\bar\tau_c$ of a field's value
and no other structural knowledge, the entropy-maximising distribution for a positive
survival time is the exponential $S_c(\dt)=e^{-\dt/\bar\tau}$, so that
\begin{equation}
  m_c(\dt) \;=\; u_c + (m_c(0)-u_c)\, e^{-\dt/\bar\tau_c},
  \label{eq:exp-m}
\end{equation}
which is \eqref{eq:surv} with the entropy-maximising survival. If the variance (or
shape) of residency is empirically available, the max-entropy form becomes the Weibull
$S_c(\dt)=e^{-(\dt/\bar\tau)^k}$--- the two-parameter generalisation that TimeLink
selects too, on generative grounds rather than entropy grounds. The ``decay rate''
$\rho_c=1/\bar\tau_c$ is therefore not a knob tuning a weight; it is the parameter of
the max-entropy survival of the field's residency distribution, estimable from labelled
history alone.

\paragraph{Step 3: estimation from data, not assumption.}
Both $m_c(0)$ and $\bar\tau_c$ are estimable (e.g.\ maximum-likelihood on the observed
gaps of matched pairs: $\bar\tau_c$ is the empirical mean residency of changed values,
and $m_c(0)$ the same-pair agreement at near-zero gap). The decay is then \emph{not
assumed to be a specific multiple of the static weight; it is inferred}. The general
form \eqref{eq:llr-time} specialises, with \eqref{eq:exp-m}, to
\begin{equation}
  w_c(\dt) = \log_2\frac{u_c + (m_c(0)-u_c)e^{-\dt/\bar\tau_c}}{u_c},
  \label{eq:weight-entropy}
\end{equation}
and the ``decay multiplier'' $\beta_c(\dt)=\log_2 m_c(\dt)/u_c$ falls from
$\log_2(m_c(0)/u_c)$ to $0$ at rate $\bar\tau_c^{-1}$, exactly as the observed
agreement distributions prescribe.

\paragraph{What this says about the TimeLink comparison.}
TimeLink is right that a survival form and a decay curve must be identified to leave the
static model. But ``ad hoc'' is not the correct characterisation of the exponential form
in the present investigation: with only a mean field lifetime available, the
\emph{maximal-entropy} survival is exponential, and the empirical-fit procedure
(mean survival time on actual matched pairs) identifies both parameters without
appeal to a specific functional form. TimeLink and the design of this project differ in
what they fit (full semi-Markov transition vs.\ per-field exponential periods), but both
place the decay inside $m$, both learn the parameters from data, and neither requires
an external knob. Where the present project goes beyond even TimeLink is
Section~\ref{sec:gap}, item~1: the fitted $\bar\tau_c$ is reported in half-lives, and
those half-lives predict (and are checked against) the measured $D_{KL}(\dt)$ curves
of each field.

\paragraph{Caveat: entropy vs.\ decision.}
Bayes factors that take entropy-compat form must still be used with a decision rule.
Entropy tells us what the evidence is (bits, sign preserve the log odds); it does not
select the match prior $\lambda$ or the threshold $\tau$, which remain decision-theoretic
choices (the costs of FP and FN). Nothing in the max-entropy derivation removes those. If anything,
it separates them cleanly: the \emph{model} is fixed by max-entropy-evidence, the
\emph{operating point} by cost.

\section{Gap analysis}
\label{sec:gap}

To the best of our knowledge:
\begin{enumerate}
  \item no published work measures per-field decay rates (half-lives) on a two-stage
        embedding-then-linkage pipeline and relates them both to blocking recall and to
        the entropy framing above;
  \item nobody has formulated the natural blocker consequence --- drop a field from a
        fused index once it has roughly exhausted its mutual information --- although
        Litoux and Ray use the idea and TimeLink keeps table lookups separate;
  \item the ``$u$ static'' assumption of the time-dependent-FS line is not carefully
        tested against \emph{cohort drift} (the value distribution itself changes over
        decades, changing $u$), a concern present in Hu et al.\ (2017), in the Li et al.\
        (2011) frequency-aware rates, and in the present project's design note;
  \item the practical weak spot remains the NC Voter file, on which temporal models are
        observed to be fragile (Christen--Gayler 2013; Hu et al.\ 2017) --- exactly where
        the project's own runs contribute additional evidence.
\end{enumerate}

\section{Notes on source status}

This survey was compiled from public records (VLDB, ACM, Springer, OpenAlex, and DBLP) and
from the accepted manuscripts of Christen--Gayler 2013 and Hu et al.\ 2017 where they could
be read in full. TimeLink (a 2026 SSRN preprint) is cited from its abstract and any
accessible body text; it was not wholly available to the survey. Where page-level facts could
not be verified in the primary text (figure numbers in Christen--Gayler 2013), the survey
says so explicitly. This review is part of the ongoing entity-resolution project, with the
purpose of positioning the ``measured per-field decay rates'' result.

\begin{thebibliography}{99}
\bibitem{fellegi1969}
I.~P.~Fellegi and A.~B.~Sunter. A theory for record linkage.
\newblock JASA 64(328):1183--1210, 1969.

\bibitem{li2011}
P.~Li, X.~L.~Dong, A.~Maurino, and D.~Srivastava. Linking temporal records.
\newblock PVLDB 4(11):956--967, 2011.

\bibitem{cg2013}
P.~Christen and R.~W.~Gayler. Adaptive temporal entity resolution on dynamic databases.
\newblock PAKDD 2013, LNAI 7819:558--569.

\bibitem{cd14}
Y.~Chiang, A.~Doan, and J.~F.~Naughton. Modeling entity evolution for temporal record
matching. \newblock SIGMOD 2014:1175--1186.

\bibitem{ll15}
F.~Li, M.~L.~Lee, W.~Hsu, and W.-C.~Tan. Linking temporal records for profiling.
\newblock SIGMOD 2015:593--605.

\bibitem{hu2017}
Y.~Hu, Q.~Wang, D.~Vatsalan, and P.~Christen. Improving temporal record linkage using
regression classification. \newblock PAKDD 2017, LNCS 10235:561--573.

\bibitem{rc2018}
T.~Ranbaduge and P.~Christen. Privacy-preserving temporal record linkage.
\newblock ICDM 2018:377--386.

\bibitem{rc2020}
T.~Ranbaduge and P.~Christen. A scalable privacy-preserving framework for temporal record
linkage. \newblock KAIS 62(1):45--78, 2020.

\bibitem{lr2026}
V.~Litoux and C.~Ray. Temporal record linkage using time decay models applied to vessel
data. \newblock MDM 2026:313--318.

\bibitem{shim2026}
K.~B.~Shim. TimeLink: time-aware record linkage with semi-Markov dynamics.
\newblock SSRN preprint 6468518, 2026.

\bibitem{winkler1993}
W.~E.~Winkler. String comparator metrics and enhanced decision rules in the
Fellegi--Sunter model of record linkage.
\newblock Proceedings of the American Statistical Association Section on Survey
Research, 1993.
\end{thebibliography}

\end{document}